\documentclass[11pt]{article}

\usepackage[margin=1in]{geometry}
\usepackage{amsmath,amssymb,amsthm}
\usepackage{array}
\usepackage{graphicx}
\usepackage[hidelinks]{hyperref}

\newtheorem{theorem}{Theorem}
\newtheorem{lemma}{Lemma}

\DeclareMathOperator{\SO}{SO}
\DeclareMathOperator{\SU}{SU}
\DeclareMathOperator{\U}{U}
\DeclareMathOperator{\so}{\mathfrak{so}}
\DeclareMathOperator{\su}{\mathfrak{su}}
\DeclareMathOperator{\rank}{rank}

\title{The Seven-Dimensional Case of the Skew-Hermitian Lie Algebra Problem}
\author{FA/Banach discovery run \texttt{fa\_banach\_001}}
\date{Candidate packet, 2026-06-13}

\begin{document}
\maketitle

\section*{Source problem}

The source paper is V. Kadets, M. Martin, and R. Paya,
``Recent progress and open questions on the numerical index of Banach
spaces'', arXiv:math/0605781. It asks for the possible dimensions of
$\mathcal Z(X)$, the Lie algebra of skew-hermitian operators, when $X$ is a
finite-dimensional real Banach space.

\begin{center}
\includegraphics[width=\textwidth]{figures/open_problem_crop.png}
\end{center}

\section*{Result}

\begin{theorem}\label{thm:seven}
Let $X$ be a seven-dimensional real Banach space. Then
\[
\dim \mathcal Z(X)\in
\{0,1,2,3,4,5,6,7,9,10,11,15,21\}.
\]
Conversely, every value in this set is attained by some seven-dimensional
real Banach space.
\end{theorem}

This is a partial result for the general Problem 27 of the source paper: it
settles the next dimension after the run's existing five- and six-dimensional
packets.

\section*{Compact representation input}

We use the following standard low-dimensional facts about compact Lie
algebras and their orthogonal representations. Here the center acts through
the skew-symmetric commutant of the semisimple part.

\begin{lemma}\label{lem:so7-table}
Let $\mathfrak g\subseteq\so(7)$ be a compact Lie algebra, let
$\mathfrak s=[\mathfrak g,\mathfrak g]$, and suppose
$\dim\mathfrak g\leq 14$. Then the following table contains all cases in
which $\mathfrak s$ can act faithfully on $\mathbb R^7$.
\begin{center}
\small
\renewcommand{\arraystretch}{1.15}
\begin{tabular}{c|c|p{0.46\linewidth}|c}
$\mathfrak s$ & $\dim\mathfrak s$
& faithful real form in dimension $\leq 7$
& max. center dim.\\ \hline
$0$ & $0$ & trivial & $3$\\
$\su(2)$ & $3$ & $3,4,5,7$-dimensional irreducibles, with trivials & $2$\\
$\su(2)\oplus\su(2)$ & $6$
& $4$-dimensional tensor form, or $3+3$, with trivials & $1$\\
$\su(2)^3$ & $9$ & $4+3$ & $0$\\
$\su(3)$ & $8$ & $\mathbb C^3$ as $\mathbb R^6$, plus at most a line & $1$\\
$\so(5)$ & $10$ & $\mathbb R^5$, plus at most two trivial dimensions & $1$\\
$\mathfrak g_2$ & $14$ & $\mathbb R^7$, the standard representation & $0$
\end{tabular}
\end{center}
In particular, no compact Lie subalgebra of $\so(7)$ has dimension $12$ or
$13$.
\end{lemma}

\begin{proof}
A compact Lie algebra is the direct sum of its center and compact simple
summands. The compact simple summands of dimension at most $14$ admitting a
non-trivial real representation of dimension at most $7$ are
\[
\su(2),\qquad \su(3),\qquad \so(5),\qquad \mathfrak g_2.
\]
Their relevant smallest faithful real representations have dimensions
$3$ or $4$ for $\su(2)$, $6$ for $\su(3)$, $5$ for $\so(5)$, and $7$ for
$\mathfrak g_2$.

The entries for sums of $\su(2)$'s follow by the same dimension count. A
faithful representation of $\su(2)^2$ of dimension at most $7$ is either the
real four-dimensional tensor form of the two fundamental representations,
with trivial summands allowed, or the sum of the two three-dimensional
adjoint representations, with at most one trivial summand. A faithful
representation of $\su(2)^3$ in dimension at most $7$ must be the
four-dimensional tensor form for two factors plus the three-dimensional
adjoint representation for the third. Four $\su(2)$ factors already require
dimension at least $8$.

The center of $\mathfrak g$ must embed into the skew-symmetric commutant of
the semisimple action. The ranks shown in the last column are immediate from
the displayed decompositions: for example, a four-dimensional
$\su(2)^2$ tensor block has no skew commutant and the largest possible
trivial complement has dimension $3$, whose skew rank is $1$; the
$\su(3)$ standard real representation has one-dimensional skew commutant
generated by the complex structure; and the $\so(5)$ standard representation
with a trivial two-plane has at most the $\so(2)$ center on that plane.
These bounds also exclude dimensions $12$ and $13$.
\end{proof}

\begin{lemma}\label{lem:sphere-transitive}
Let $X$ be a finite-dimensional real Banach space and suppose
$\mathcal Z(X)$ contains a compact semisimple algebra acting on an invariant
Euclidean block in one of the following standard ways:
\[
\su(3)\curvearrowright \mathbb C^3\cong\mathbb R^6,
\qquad
\mathfrak g_2\curvearrowright\mathbb R^7.
\]
Then the corresponding block is actually invariant under the full orthogonal
group of that block.
\end{lemma}

\begin{proof}
The connected groups $\SU(3)$ on $\mathbb C^3$ and $G_2$ on $\mathbb R^7$
are transitive on the Euclidean unit spheres of their standard
representations. Hence a norm invariant under either group depends on the
block variable only through its Euclidean length, with any remaining
coordinates held fixed. Therefore the norm is invariant under the full
orthogonal group of that block.
\end{proof}

\section*{Proof of Theorem \ref{thm:seven}}

The source paper recalls two facts. First, by Auerbach's theorem one may put
an inner product on $X$ so that $\mathcal Z(X)$ is a compact Lie subalgebra
of $\so(7)$. Second, Rosenthal's gap theorem says that for
$\dim X=n$:
\[
\dim\mathcal Z(X)>\frac{(n-1)(n-2)}2
\quad\Longrightarrow\quad
X\hbox{ is Hilbert,}
\]
so the dimension is $n(n-1)/2$; and equality at
$(n-1)(n-2)/2$ occurs exactly for a non-Euclidean absolute sum of
$\mathbb R$ and a Hilbert space of dimension $n-1$.

For $n=7$, this leaves only dimensions at most $15$, together with the
Hilbert value $21$. The equality case $15$ is possible and already described
by Rosenthal. Thus it remains to rule out the values
\[
8,\quad 12,\quad 13,\quad 14.
\]
Lemma \ref{lem:so7-table} rules out $12$ and $13$ even as dimensions of
compact Lie subalgebras of $\so(7)$.

If $\dim\mathcal Z(X)=8$, Lemma \ref{lem:so7-table} forces the semisimple
part to be $\su(3)$ acting on a six-dimensional block $\mathbb C^3$, with at
most one fixed real line. By Lemma \ref{lem:sphere-transitive}, the norm is
then invariant under $\SO(6)$ on that block, so $\mathcal Z(X)$ contains
$\so(6)$ and has dimension at least $15$, a contradiction.

If $\dim\mathcal Z(X)=14$, Lemma \ref{lem:so7-table} forces the standard
$\mathfrak g_2$ action on $\mathbb R^7$. Since $G_2$ is transitive on the
unit sphere in $\mathbb R^7$, Lemma \ref{lem:sphere-transitive} implies that
the norm is Euclidean and $\mathcal Z(X)=\so(7)$, of dimension $21$. This is
again a contradiction.

So the only possible dimensions are
\[
\{0,1,2,3,4,5,6,7,9,10,11,15,21\}.
\]

It remains to realize them. Let $H_k$ denote the $k$-dimensional real Hilbert
space. Finite $\ell_\infty$-sums of Hilbert blocks realize all the listed
values, because the connected isometry group is the product of the orthogonal
groups on the Hilbert blocks and
$\dim\so(k)=k(k-1)/2$:
\[
\begin{array}{c|l}
\dim\mathcal Z(X) & X\\ \hline
0 & \ell_\infty^7\\
1 & H_2\oplus_\infty \ell_\infty^5\\
2 & H_2\oplus_\infty H_2\oplus_\infty \ell_\infty^3\\
3 & H_3\oplus_\infty \ell_\infty^4\\
4 & H_3\oplus_\infty H_2\oplus_\infty \ell_\infty^2\\
5 & H_3\oplus_\infty H_2\oplus_\infty H_2\\
6 & H_4\oplus_\infty \ell_\infty^3\\
7 & H_4\oplus_\infty H_2\oplus_\infty \mathbb R\\
9 & H_4\oplus_\infty H_3\\
10 & H_5\oplus_\infty \ell_\infty^2\\
11 & H_5\oplus_\infty H_2\\
15 & H_6\oplus_\infty \mathbb R\\
21 & H_7.
\end{array}
\]
This proves the classification.

\section*{Verification notes}

The proof has three review points:
\begin{enumerate}
\item Rosenthal's gap theorem is used exactly as quoted in the source paper.
\item Lemma \ref{lem:so7-table} is the finite compact-representation check.
It is the main human-verification target.
\item The exclusions of $8$ and $14$ use only the standard sphere
transitivity of $\SU(3)$ on $S^5\subset\mathbb C^3$ and $G_2$ on
$S^6\subset\mathbb R^7$.
\end{enumerate}
The script \texttt{code/check\_seven\_dimension\_values.py} confirms the
Euclidean-block construction values and the four missing values below the
Rosenthal threshold. It is a sanity check, not a proof of the representation
lemma.

\section*{Novelty check}

Bounded local searches checked the run registry, solution index, attempts,
and parsed source for \texttt{0605781}, \texttt{dim Z(X)},
\texttt{skew-hermitian}, and the existing five- and six-dimensional packets.
The run contains no previous seven-dimensional classification.

Bounded web searches on 2026-06-13 used exact or close phrases involving:
\begin{itemize}
\item \texttt{dim Z(X)}, \texttt{7-dimensional}, and \texttt{Banach space};
\item \texttt{skew-hermitian}, \texttt{possible values}, and
\texttt{dimension Z(X)};
\item Rosenthal's phrase \texttt{The Lie algebra of a Banach space}.
\end{itemize}
They found the source arXiv page and adjacent background, but no direct later
classification of the seven-dimensional case.

\section*{Scope}

This packet settles only the seven-dimensional subcase of Problem 27. It does
not claim the general all-dimensional classification.

\end{document}
